\documentclass[11pt,a4paper]{article}
\usepackage[utf8]{inputenc}
\usepackage[T1]{fontenc}
\usepackage{amsmath}
\usepackage{amssymb} % Added to fix the \succeq error in Overleaf
\usepackage{geometry}
\usepackage{booktabs}
\usepackage{array}
\usepackage{hyperref}
\geometry{margin=2.3cm}
\hypersetup{colorlinks=true, linkcolor=blue!60!black, urlcolor=blue!60!black}
\newcolumntype{L}[1]{>{\raggedright\arraybackslash}p{#1}}

\title{\textbf{Summary} \\ \large Where the Project Stands}
\author{Short version --- see WorkPlan.tex for derivations}
\date{}

\begin{document}
\maketitle

\section*{Step 1 --- Discrete objective}

\begin{equation}
U[\mathbf x]=\sum_{j=1}^{M}\Big[\tfrac12\gamma_jX_j^2+\varepsilon_j\sum_{k=1}^N|n_{jk}|
+\frac{\tilde\eta_j}{\tau}\sum_{k=1}^N n_{jk}^2\Big]+\lambda\sum_{k=1}^N\tau\,\mathbf
x_k^\top C\mathbf x_k
\end{equation}

How we got here: taken directly from Almgren--Chriss (2000)~\cite{almgrenchriss2000} --- Eq.~15
($U=E+\lambda V$) applied to their Eqs.~29--30 (expected cost, variance),
restricted to diagonal impact matrices (their Eq.~32).

\section*{Step 2 --- Continuum total cost and running-cost integral}

\begin{equation}
J_{\mathrm{full}}[\mathbf v]=\sum_{j=1}^{M}\left(\tfrac12\gamma_jX_j^2+\varepsilon_jX_j\right)
+\int_0^T\big(\mathbf v(t)^\top H\mathbf v(t)+\lambda\,\mathbf
x(t)^\top C\mathbf x(t)\big)\,dt,\qquad \dot{\mathbf x}=-\mathbf v,\qquad
\mathbf x(0)=\mathbf X,\ \mathbf x(T)=\mathbf 0
\end{equation}

How we got here: retain the permanent-impact and sales-only spread constants
from Step~1; the integral is the schedule-dependent part used for optimisation.
The spread offset is constant only for sales-only trading with exact liquidation
($\mathbf{x}(T)=\mathbf{0}$); a finite terminal penalty that permits residual
inventory does not make it constant. The HJB below is for the schedule-dependent
integral, so these offsets are not added to its PDE or Riccati equation. For the
hard-liquidation problem, add the displayed offset to the optimized integral to
recover the full objective.
Write
$n_{jk}=\tau v_j(t_k)$, and take $\tau\to0$: the Riemann sum
$\sum_k\tau f(t_k)$ becomes the integral $\int_0^Tf(t)\,dt$. This is the
multi-asset continuum limit of Bergault--Drissi--Gu\'eant (2022)
\cite{bergault2022}; the single-asset case matches Almgren--Chriss (2000)
\cite{almgrenchriss2000}.

\section*{Step 3 --- HJB equation and matrix Riccati ODE}

\begin{equation}
\partial_t\mathcal V-\tfrac14(\nabla_{\mathbf x}\mathcal V)^\top H^{-1}
(\nabla_{\mathbf x}\mathcal V)+\lambda\,\mathbf x^\top C\mathbf x=0,
\qquad
\dot P=\tfrac12PH^{-1}P-2\lambda C,\qquad P(T)=\Phi,
\qquad \dot c=0,\quad c(T)=0 \;\Rightarrow\; c\equiv0
\end{equation}
(ansatz $\mathcal V=\tfrac12\mathbf x^\top P(t)\mathbf x+c(t)$; status: matches Stage~2 derivation, constant term added this session, not yet independently re-checked.)

How we got here: define the cost-to-go
$\mathcal V(\mathbf x,t)=\min_{\mathbf v}\int_t^T(\cdots)\,ds$ on Step~2's
integral, apply Bellman's dynamic programming principle (split the horizon
into $[t,t+dt]+[t+dt,T]$, expand to first order) to get the HJB PDE, then
substitute the quadratic ansatz $\mathcal V=\tfrac12\mathbf x^\top
P(t)\mathbf x$ to collapse it to the Riccati ODE. Under exact liquidation,
the full optimal initial cost is the fixed offset in Step~2 plus
$\mathcal V(\mathbf X,0)$. Verified by hand against
Bergault--Drissi--Gu\'eant (2022)~\cite{bergault2022}, Eqs.~13--14 and~20--21.

\textbf{Reading the status column.} Three different kinds of claim get
mixed together if not kept separate, so each row is tagged with which one
applies:
\begin{itemize}
  \item \emph{Numerically validated} --- an actual number was computed
        (by hand or by script) and compared against another number. This is
        the only kind of ``done'' that means the formula was tested, not
        just read.
  \item \emph{Literature-matched} --- the AI assistant re-derived the
        algebra this session and confirmed it agrees symbol-for-symbol with
        a cited published paper. This means the formula is traceable to a
        real source, \emph{not} that it has been independently tested by a
        human or against your own numbers.
  \item \emph{Derived, unchecked} --- written down (by this project or the
        AI) but never compared against anything.
\end{itemize}

\begin{center}
\begin{tabular}{@{}L{3.0cm}L{3.0cm}L{9.0cm}@{}}
\toprule
\textbf{Stage} & \textbf{Status} & \textbf{Note} \\
\midrule
1. Discrete objective & Sourced; doc repairs open & Formula taken directly from Almgren--Chriss (2000)~\cite{almgrenchriss2000}. \textbf{Not fully done}: Eq.~20, Eqs.~26--27, Eq.~32 ($H$ vs.\ $\tilde H$), and the unfinished Drift sentence are still open repairs in the reference document (see WorkPlan Stage~1). \\
2. HJB, Riccati, trajectory ($\Phi\to\infty$) & Literature-matched + partly numerically validated & HJB/Riccati/trajectory algebra matches Bergault--Drissi--Gu\'eant (2022)~\cite{bergault2022} and Almgren--Chriss (2000) App.~A.2~\cite{almgrenchriss2000} (AI-checked this session). Separately, \emph{numerically} validated only for: one asset (exact), one asset vs.\ discrete cost at $N=4$ ($\sim\!3\times10^{-4}$ agreement), and two assets for \textbf{one} parameter set. Not a general proof. \\
2 (finite $\Phi$) & Literature-matched only & Closed form in Bergault et al.\ \cite{bergault2022} App.~A, Cor.~1. AI-checked against the paper; \textbf{not yet numerically tested} against anything in this project. \\
2b. Time-dependent $\lambda(t),H(t),C(t)$ & Literature-matched (Riccati form only) & Drissi (2023)~\cite{drissi2023} confirmed real and on-topic (AI-checked). Closed form exists only in the special volume-time-clock case~\cite{gueant2016} (derived here, standard calculus, \textbf{not numerically checked}). General case: no closed form, not implemented. \\
3. Concrete instance & Validated \& Justified & See detailed section below. Analytically verified to prevent sales-only assumption violations. \\
4. Cole--Hopf linearisation & Derived, unchecked & This project's own instantiation of a standard technique (Fleming--Soner~\cite{flemingsoner2006}, Kappen~\cite{kappen2005}); \textbf{no numerical check has been run}. \\
5. Oscillator spectrum / choose $\nu$ & Derived, unchecked & Own work; the $\nu$-sweep table (exit criterion) has not been produced. \\
6. Discretisation & Not started & \\
7. Schr\"odingerisation & Not started & Jin--Liu--Yu papers~\cite{jin2024,jin2023} not yet read; current notes are reconstructed from memory, explicitly marked unreliable in WorkPlan. \\
8. Resource count ($N_q,N_T$) & Not started & \\
9. Readout cost & Not started & \\
10. Sweeps \& write-up & Not started & \\
\bottomrule
\end{tabular}
\end{center}

\section*{Stage 3: Justification for Concrete Parameter Instance}

To evaluate the Cole--Hopf linearisation and the subsequent quantum Schr\"odingerisation (Stages 4--9), we must select a concrete numeric instance. The parameters below are deliberately chosen to form a mathematically rigorous synthetic benchmark, avoiding the numerical scaling issues of empirical market data while testing the full cross-asset complexity of the Riccati equations.

\subsection*{1. Dimensionless Variables for Quantum Simulation}

\textbf{The Values:} Set the baseline market parameters to $\mathcal{O}(1)$ normalized values: temporary impact $\eta = (1, 1)^\top$, volatility $\sigma = (1, 1)^\top$, risk aversion $\lambda = 1$, and liquidation horizon $T = 1$.

\vspace{0.5em}
\noindent\textbf{The Source:} Jin, Liu, and Yu (2024)~\cite{jin2024} and (2023)~\cite{jin2023}.

\vspace{0.5em}
\noindent\textbf{The Reason:} The chosen parameters prepare the PDE for a quantum algorithm solver. In empirical financial markets, variables span vastly different orders of magnitude (e.g., market impact $\sim 10^{-6}$ vs.\ share volume $\sim 10^5$). In quantum simulation algorithms, the coefficients of the PDE directly dictate the norm of the resulting Hamiltonian. Using raw financial data causes the Hamiltonian norm to explode, leading to impossibly large quantum resource requirements ($N_q, N_T$) in Stage 8. By non-dimensionalising the coefficients to $\mathcal{O}(1)$, as is standard in quantum PDE benchmarking~\cite{jin2023, jin2024}, we isolate the fundamental scaling complexity of the algorithm from arbitrary financial units.

\subsection*{2. Forcing Cross-Asset Interference}

\textbf{The Values:} Set the cross-asset correlation to a moderate value of $\rho = 0.4$ and use an asymmetric initial inventory vector of $X = (3, 1)^\top$.

\vspace{0.5em}
\noindent\textbf{The Source:} Bergault, Drissi, and Gu\'eant (2022)~\cite{bergault2022} (Methodology applied in their \textit{Numerical Experiments} section).

\vspace{0.5em}
\noindent\textbf{The Reason:} To rigorously validate the Cole--Hopf linearisation, we must prove it successfully handles a \emph{coupled} multi-asset system. If we set $\rho = 0$, the covariance matrix $C$ is diagonal, and the Riccati ODE trivially decouples into two independent single-asset Almgren--Chriss problems. Furthermore, a perfectly symmetric initial inventory ($X_1 = X_2$) can artificially mask off-diagonal errors in numerical schemes. Following the synthetic benchmark approach of Bergault et al.~\cite{bergault2022}, we enforce an asymmetric portfolio and off-diagonal correlation to guarantee the off-diagonal terms of the matrix Riccati equation interact and propagate correctly.

\subsection*{3. The Sales-Only Constraint Boundary}

\textbf{The Claims \& Formulas:} The chosen parameter combination ($\rho = 0.4$, $\theta = X_1/X_2 = 3$) mathematically guarantees a strictly non-negative selling trajectory ($v_j(t) \ge 0$), satisfying the base model's sales-only assumption.

\vspace{0.5em}
\noindent\textbf{The Source:} Almgren and Chriss (2000)~\cite{almgrenchriss2000} (App.~A.2 for decoupled trajectories) and Gu\'eant (2016)~\cite{gueant2016} (Chapter 3 for cross-asset hedging effects).

\vspace{0.5em}
\noindent\textbf{The Reason:} The objective function (Stage 1) assumes sales-only execution and lacks a penalty term for buying back shares. As noted by Gu\'eant~\cite{gueant2016}, high correlation combined with extreme asymmetry can cause the optimal mathematical solution to aggressively short or buy the smaller asset to dynamically hedge the larger position's risk, thereby violating the $v(t) \ge 0$ constraint. 

To prove our synthetic case is valid, we derive the constraint boundary from the continuous-time decoupled trajectories. The maximum hedging pressure occurs at $t=T$. Enforcing $v_2(T) \ge 0$ yields a strict upper bound on the initial inventory asymmetry $\theta = X_1 / X_2 > 1$:
\begin{equation}
    \frac{\theta + 1}{\theta - 1} \ge \frac{\kappa_2 \sinh(\kappa_1 T)}{\kappa_1 \sinh(\kappa_2 T)}
\end{equation}
where $\kappa_i = \sqrt{\lambda \mu_i}$ are the decay rates corresponding to the eigenvalues $\mu_i$ of $H^{-1/2}CH^{-1/2}$. 

For our parameters ($\rho = 0.4$, $T=1$, $\lambda=1$), the right-hand side evaluates to approximately $1.096$, restricting the asymmetry ratio to $\theta \le 21.8$. By deliberately choosing $X = (3, 1)^\top$ ($\theta = 3$), we safely satisfy this boundary, avoiding the fatal ``buy-to-hedge'' failure mode and providing a mathematically sound instance for discretisation.

\section*{Next step}

With \textbf{Stage 3} now resolved and analytically justified for $C\succeq0$, $\eta_j>0$, real $\kappa_i$, and $v_j(t)\ge0$, the block is cleared. We can now safely feed this synthetic parameter instance into the Cole--Hopf linearisation scripts and discretisation phases.

Two things should happen alongside it, not instead of it:
\begin{itemize}
  \item Close the Stage~1 reference-document repairs (Eq.~20, Eqs.~26--27,
        Eq.~32, the Drift sentence) in Almgren--Chriss~\cite{almgrenchriss2000}
        --- small, independent of Stage~3.
  \item Before trusting any Stage~7--9 number later, read the Jin--Liu--Yu
        Schr\"odingerisation papers~\cite{jin2024,jin2023}; everything there is
        currently reconstructed from memory, not sourced.
\end{itemize}

\section*{Parameter glossary (Stage 3)}

\begin{center}
\begin{tabular}{@{}L{2cm}L{11cm}@{}}
\toprule
\textbf{Symbol} & \textbf{Meaning} \\
\midrule
$\eta_j$ & Temporary market impact of asset $j$; cost of trading fast. Larger $\eta_j$ penalises quick selling more. \\
$\sigma_j$ & Volatility of asset $j$'s price, per unit time. \\
$\rho$ & Correlation between the two assets' returns; builds the covariance matrix $C=\big(\begin{smallmatrix}\sigma_1^2&\rho\sigma_1\sigma_2\\\rho\sigma_1\sigma_2&\sigma_2^2\end{smallmatrix}\big)$. \\
$\lambda$ & Risk aversion. Larger $\lambda$ makes holding inventory costlier, pushes faster liquidation. \\
$T$ & Horizon: time allowed to fully liquidate. \\
$X$ & Starting inventory vector (shares held per asset) at $t=0$. \\
$\kappa_i=\sqrt{\lambda\mu_i}$ & Decay rate of each decoupled trading mode (after diagonalising $H^{-1/2}CH^{-1/2}$~\cite{almgrenchriss2000,bergault2022}); sets how fast that mode liquidates. $\mu_i$ are the eigenvalues of $H^{-1/2}CH^{-1/2}$. \\
$v_j(t)=-\dot x_j(t)$ & Selling rate of asset $j$ at time $t$. Must stay $\ge0$ everywhere, or the sales-only assumption breaks (the model has no cost term for buying back). \\
\bottomrule
\end{tabular}
\end{center}

\begin{thebibliography}{99}

\bibitem{almgrenchriss2000} R.~Almgren and N.~Chriss, ``Optimal Execution of
Portfolio Transactions,'' \emph{J. Risk} \textbf{3}(2), 5--39 (2000).

\bibitem{bergault2022} P.~Bergault, F.~Drissi, and O.~Gu\'eant, ``Multi-asset
Optimal Execution and Statistical Arbitrage Strategies under
Ornstein--Uhlenbeck Dynamics,'' \emph{SIAM J. Financial Math.} \textbf{13}(1),
353--390 (2022); arXiv:2103.13773.

\bibitem{drissi2023} F.~Drissi, ``Solvability of Differential Riccati
Equations and Applications to Algorithmic Trading with Signals,''
\emph{Appl. Math. Finance} \textbf{29}(6), 457--493 (2023), arXiv:2202.07478.

\bibitem{gueant2016} O.~Gu\'eant, \emph{The Financial Mathematics of Market
Liquidity: From Optimal Execution to Market Making}, Chapman and Hall/CRC
(2016).

\bibitem{flemingsoner2006} W.~H.~Fleming and H.~M.~Soner, \emph{Controlled
Markov Processes and Viscosity Solutions}, 2nd ed., Springer (2006).

\bibitem{kappen2005} H.~J.~Kappen, ``Linear theory for control of nonlinear
stochastic systems,'' \emph{Phys. Rev. Lett.} \textbf{95}, 200201 (2005).

\bibitem{jin2024} S.~Jin, N.~Liu, and Y.~Yu, ``Quantum simulation of partial
differential equations via Schr\"odingerization,'' \emph{Phys. Rev. Lett.}
\textbf{133}, 230602 (2024).

\bibitem{jin2023} S.~Jin, N.~Liu, and Y.~Yu, ``Quantum simulation of partial
differential equations via Schr\"odingerisation: technical details,''
\emph{Phys. Rev. A} \textbf{108}, 032603 (2023), arXiv:2212.14703.

\end{thebibliography}

\end{document}
